% Seminar 3 - Pagina 7. Fragment pentru
% % Seminar 3 - Pagina 7. Fragment pentru
% % Seminar 3 - Pagina 7. Fragment pentru
% % Seminar 3 - Pagina 7. Fragment pentru
% \input{seminars/pages/seminar3-pagina7-ro}; fără preambul sau pagină nouă.
% NOTE EDITORIALE:
% - Pagina continuă exemplele din paginile 5-6. Numerotarea vizibilă
%   (ii), (iii), (iv) este păstrată, chiar dacă diferă pe paginile vecine.
% - Cele două diagrame triunghiulare sunt reproduse cu tikz-cd.
% - Pentru aplicația argumentului pe C*, manuscrisul nu precizează alegerea
%   argumentului; normalizarea z/|z| exprimă corespondența indicată.
% - La reflexivitatea din (iv), manuscrisul pare să scrie „0 in R”, nu
%   „0 in Q”, cum ar cere demonstrația. Textul aparent este păstrat.
% - Diagrama cu o aplicație R -> R/~ și condiția f(x)=f(y) iff x-y in Q
%   sunt barate în manuscris; nu sunt prezentate drept afirmații valide.
% - R/~ ~ R exprimă o echivalență de cardinalitate, nu o bijecție canonică.
% - Propoziția de după p:R -> R/~ pare să spună că R nu este, ca atare,
%   un sistem de reprezentanți; descifrarea exactă rămâne nesigură.

\subsection*{Seminar 3 - Pagina 7}
\begingroup
\setlength{\abovedisplayskip}{5pt}
\setlength{\belowdisplayskip}{5pt}
\setlength{\abovedisplayshortskip}{3pt}
\setlength{\belowdisplayshortskip}{3pt}

\noindent\textbf{(ii) Clasele numerelor cu același modul (continuare).}
Proiecția canonică $p$ și aplicația modulului $f$ dau diagrama
\[
\begin{tikzcd}[column sep=large,row sep=small]
 \mathbb C\arrow[r,"p"]\arrow[d,"f"'] & \mathbb C/{\sim}\\
 {[0,+\infty)}\arrow[ur,dashed,"\overline f"'] &
\end{tikzcd}
\qquad f(z)=|z|.
\]
Aici $\sim=R_f$, adică
\[
 x\sim y\iff f(x)=f(y).
\]

\noindent\textbf{(iii) Clasele numerelor cu același argument (continuare).}
Pentru $\mathbb C^*=\mathbb C\setminus\{0\}$ și cercul unitate $S'$:
\[
\begin{tikzcd}[column sep=large,row sep=small]
 \mathbb C^*\arrow[r,"p"]\arrow[d,"f"'] & \mathbb C^*/{\sim}\\
 S'\arrow[ur,dashed,"\overline f"'] &
\end{tikzcd}
\]
Scriind $z=r(\cos t+i\sin t)$, în notițe apare
\[
 f(z)=f\bigl(r(\cos t+i\sin t)\bigr)=\cos t+i\sin t.
\]

\noindent\textbf{(iv) Echivalența modulo $\mathbb Q$ pe $\mathbb R$.}
Definim
\[
 x\sim y\overset{\mathrm{def}}\iff x-y\in\mathbb Q,
 \qquad \mathbb R/{\sim}\cong\mathbb R.
\]
Verificăm cele trei proprietăți:
\[
\begin{aligned}
\text{(R)}\quad &x\sim x\quad(0\in\mathbb R),\\
\text{(S)}\quad &x-y\in\mathbb Q\Longrightarrow y-x\in\mathbb Q,\\
\text{(T)}\quad &x-y\in\mathbb Q,\ y-z\in\mathbb Q
      \Longrightarrow x-z\in\mathbb Q.
\end{aligned}
\]
Proiecția este $p:\mathbb R\to\mathbb R/{\sim}$.
Notițele remarcă faptul că $\mathbb R$ nu este, ca atare,
un sistem de reprezentanți. Pentru un sistem de reprezentanți
$S=(x_i)_{i\in I}$,
\[
 p|_S:S\longrightarrow\mathbb R/{\sim}\quad\text{este bijecție}.
\]
Deoarece
\[
 \mathbb R=\bigcup_{i\in I}\widehat{x_i},\qquad
 \widehat{x_i}=\{y\in\mathbb R\mid y-x_i\in\mathbb Q\}
   =x_i+\mathbb Q\cong\mathbb Q,
\]
fiecare $\widehat{x_i}$ este o mulțime numărabilă. Dacă $I$ ar fi
numărabilă, atunci $\mathbb R$ ar fi numărabilă - contradicție.
Prin urmare, $I$ este nenumărabilă.

\endgroup
; fără preambul sau pagină nouă.
% NOTE EDITORIALE:
% - Pagina continuă exemplele din paginile 5-6. Numerotarea vizibilă
%   (ii), (iii), (iv) este păstrată, chiar dacă diferă pe paginile vecine.
% - Cele două diagrame triunghiulare sunt reproduse cu tikz-cd.
% - Pentru aplicația argumentului pe C*, manuscrisul nu precizează alegerea
%   argumentului; normalizarea z/|z| exprimă corespondența indicată.
% - La reflexivitatea din (iv), manuscrisul pare să scrie „0 in R”, nu
%   „0 in Q”, cum ar cere demonstrația. Textul aparent este păstrat.
% - Diagrama cu o aplicație R -> R/~ și condiția f(x)=f(y) iff x-y in Q
%   sunt barate în manuscris; nu sunt prezentate drept afirmații valide.
% - R/~ ~ R exprimă o echivalență de cardinalitate, nu o bijecție canonică.
% - Propoziția de după p:R -> R/~ pare să spună că R nu este, ca atare,
%   un sistem de reprezentanți; descifrarea exactă rămâne nesigură.

\subsection*{Seminar 3 - Pagina 7}
\begingroup
\setlength{\abovedisplayskip}{5pt}
\setlength{\belowdisplayskip}{5pt}
\setlength{\abovedisplayshortskip}{3pt}
\setlength{\belowdisplayshortskip}{3pt}

\noindent\textbf{(ii) Clasele numerelor cu același modul (continuare).}
Proiecția canonică $p$ și aplicația modulului $f$ dau diagrama
\[
\begin{tikzcd}[column sep=large,row sep=small]
 \mathbb C\arrow[r,"p"]\arrow[d,"f"'] & \mathbb C/{\sim}\\
 {[0,+\infty)}\arrow[ur,dashed,"\overline f"'] &
\end{tikzcd}
\qquad f(z)=|z|.
\]
Aici $\sim=R_f$, adică
\[
 x\sim y\iff f(x)=f(y).
\]

\noindent\textbf{(iii) Clasele numerelor cu același argument (continuare).}
Pentru $\mathbb C^*=\mathbb C\setminus\{0\}$ și cercul unitate $S'$:
\[
\begin{tikzcd}[column sep=large,row sep=small]
 \mathbb C^*\arrow[r,"p"]\arrow[d,"f"'] & \mathbb C^*/{\sim}\\
 S'\arrow[ur,dashed,"\overline f"'] &
\end{tikzcd}
\]
Scriind $z=r(\cos t+i\sin t)$, în notițe apare
\[
 f(z)=f\bigl(r(\cos t+i\sin t)\bigr)=\cos t+i\sin t.
\]

\noindent\textbf{(iv) Echivalența modulo $\mathbb Q$ pe $\mathbb R$.}
Definim
\[
 x\sim y\overset{\mathrm{def}}\iff x-y\in\mathbb Q,
 \qquad \mathbb R/{\sim}\cong\mathbb R.
\]
Verificăm cele trei proprietăți:
\[
\begin{aligned}
\text{(R)}\quad &x\sim x\quad(0\in\mathbb R),\\
\text{(S)}\quad &x-y\in\mathbb Q\Longrightarrow y-x\in\mathbb Q,\\
\text{(T)}\quad &x-y\in\mathbb Q,\ y-z\in\mathbb Q
      \Longrightarrow x-z\in\mathbb Q.
\end{aligned}
\]
Proiecția este $p:\mathbb R\to\mathbb R/{\sim}$.
Notițele remarcă faptul că $\mathbb R$ nu este, ca atare,
un sistem de reprezentanți. Pentru un sistem de reprezentanți
$S=(x_i)_{i\in I}$,
\[
 p|_S:S\longrightarrow\mathbb R/{\sim}\quad\text{este bijecție}.
\]
Deoarece
\[
 \mathbb R=\bigcup_{i\in I}\widehat{x_i},\qquad
 \widehat{x_i}=\{y\in\mathbb R\mid y-x_i\in\mathbb Q\}
   =x_i+\mathbb Q\cong\mathbb Q,
\]
fiecare $\widehat{x_i}$ este o mulțime numărabilă. Dacă $I$ ar fi
numărabilă, atunci $\mathbb R$ ar fi numărabilă - contradicție.
Prin urmare, $I$ este nenumărabilă.

\endgroup
; fără preambul sau pagină nouă.
% NOTE EDITORIALE:
% - Pagina continuă exemplele din paginile 5-6. Numerotarea vizibilă
%   (ii), (iii), (iv) este păstrată, chiar dacă diferă pe paginile vecine.
% - Cele două diagrame triunghiulare sunt reproduse cu tikz-cd.
% - Pentru aplicația argumentului pe C*, manuscrisul nu precizează alegerea
%   argumentului; normalizarea z/|z| exprimă corespondența indicată.
% - La reflexivitatea din (iv), manuscrisul pare să scrie „0 in R”, nu
%   „0 in Q”, cum ar cere demonstrația. Textul aparent este păstrat.
% - Diagrama cu o aplicație R -> R/~ și condiția f(x)=f(y) iff x-y in Q
%   sunt barate în manuscris; nu sunt prezentate drept afirmații valide.
% - R/~ ~ R exprimă o echivalență de cardinalitate, nu o bijecție canonică.
% - Propoziția de după p:R -> R/~ pare să spună că R nu este, ca atare,
%   un sistem de reprezentanți; descifrarea exactă rămâne nesigură.

\subsection*{Seminar 3 - Pagina 7}
\begingroup
\setlength{\abovedisplayskip}{5pt}
\setlength{\belowdisplayskip}{5pt}
\setlength{\abovedisplayshortskip}{3pt}
\setlength{\belowdisplayshortskip}{3pt}

\noindent\textbf{(ii) Clasele numerelor cu același modul (continuare).}
Proiecția canonică $p$ și aplicația modulului $f$ dau diagrama
\[
\begin{tikzcd}[column sep=large,row sep=small]
 \mathbb C\arrow[r,"p"]\arrow[d,"f"'] & \mathbb C/{\sim}\\
 {[0,+\infty)}\arrow[ur,dashed,"\overline f"'] &
\end{tikzcd}
\qquad f(z)=|z|.
\]
Aici $\sim=R_f$, adică
\[
 x\sim y\iff f(x)=f(y).
\]

\noindent\textbf{(iii) Clasele numerelor cu același argument (continuare).}
Pentru $\mathbb C^*=\mathbb C\setminus\{0\}$ și cercul unitate $S'$:
\[
\begin{tikzcd}[column sep=large,row sep=small]
 \mathbb C^*\arrow[r,"p"]\arrow[d,"f"'] & \mathbb C^*/{\sim}\\
 S'\arrow[ur,dashed,"\overline f"'] &
\end{tikzcd}
\]
Scriind $z=r(\cos t+i\sin t)$, în notițe apare
\[
 f(z)=f\bigl(r(\cos t+i\sin t)\bigr)=\cos t+i\sin t.
\]

\noindent\textbf{(iv) Echivalența modulo $\mathbb Q$ pe $\mathbb R$.}
Definim
\[
 x\sim y\overset{\mathrm{def}}\iff x-y\in\mathbb Q,
 \qquad \mathbb R/{\sim}\cong\mathbb R.
\]
Verificăm cele trei proprietăți:
\[
\begin{aligned}
\text{(R)}\quad &x\sim x\quad(0\in\mathbb R),\\
\text{(S)}\quad &x-y\in\mathbb Q\Longrightarrow y-x\in\mathbb Q,\\
\text{(T)}\quad &x-y\in\mathbb Q,\ y-z\in\mathbb Q
      \Longrightarrow x-z\in\mathbb Q.
\end{aligned}
\]
Proiecția este $p:\mathbb R\to\mathbb R/{\sim}$.
Notițele remarcă faptul că $\mathbb R$ nu este, ca atare,
un sistem de reprezentanți. Pentru un sistem de reprezentanți
$S=(x_i)_{i\in I}$,
\[
 p|_S:S\longrightarrow\mathbb R/{\sim}\quad\text{este bijecție}.
\]
Deoarece
\[
 \mathbb R=\bigcup_{i\in I}\widehat{x_i},\qquad
 \widehat{x_i}=\{y\in\mathbb R\mid y-x_i\in\mathbb Q\}
   =x_i+\mathbb Q\cong\mathbb Q,
\]
fiecare $\widehat{x_i}$ este o mulțime numărabilă. Dacă $I$ ar fi
numărabilă, atunci $\mathbb R$ ar fi numărabilă - contradicție.
Prin urmare, $I$ este nenumărabilă.

\endgroup
; fără preambul sau pagină nouă.
% NOTE EDITORIALE:
% - Pagina continuă exemplele din paginile 5-6. Numerotarea vizibilă
%   (ii), (iii), (iv) este păstrată, chiar dacă diferă pe paginile vecine.
% - Cele două diagrame triunghiulare sunt reproduse cu tikz-cd.
% - Pentru aplicația argumentului pe C*, manuscrisul nu precizează alegerea
%   argumentului; normalizarea z/|z| exprimă corespondența indicată.
% - La reflexivitatea din (iv), manuscrisul pare să scrie „0 in R”, nu
%   „0 in Q”, cum ar cere demonstrația. Textul aparent este păstrat.
% - Diagrama cu o aplicație R -> R/~ și condiția f(x)=f(y) iff x-y in Q
%   sunt barate în manuscris; nu sunt prezentate drept afirmații valide.
% - R/~ ~ R exprimă o echivalență de cardinalitate, nu o bijecție canonică.
% - Propoziția de după p:R -> R/~ pare să spună că R nu este, ca atare,
%   un sistem de reprezentanți; descifrarea exactă rămâne nesigură.

\subsection*{Seminar 3 - Pagina 7}
\begingroup
\setlength{\abovedisplayskip}{5pt}
\setlength{\belowdisplayskip}{5pt}
\setlength{\abovedisplayshortskip}{3pt}
\setlength{\belowdisplayshortskip}{3pt}

\noindent\textbf{(ii) Clasele numerelor cu același modul (continuare).}
Proiecția canonică $p$ și aplicația modulului $f$ dau diagrama
\[
\begin{tikzcd}[column sep=large,row sep=small]
 \mathbb C\arrow[r,"p"]\arrow[d,"f"'] & \mathbb C/{\sim}\\
 {[0,+\infty)}\arrow[ur,dashed,"\overline f"'] &
\end{tikzcd}
\qquad f(z)=|z|.
\]
Aici $\sim=R_f$, adică
\[
 x\sim y\iff f(x)=f(y).
\]

\noindent\textbf{(iii) Clasele numerelor cu același argument (continuare).}
Pentru $\mathbb C^*=\mathbb C\setminus\{0\}$ și cercul unitate $S'$:
\[
\begin{tikzcd}[column sep=large,row sep=small]
 \mathbb C^*\arrow[r,"p"]\arrow[d,"f"'] & \mathbb C^*/{\sim}\\
 S'\arrow[ur,dashed,"\overline f"'] &
\end{tikzcd}
\]
Scriind $z=r(\cos t+i\sin t)$, în notițe apare
\[
 f(z)=f\bigl(r(\cos t+i\sin t)\bigr)=\cos t+i\sin t.
\]

\noindent\textbf{(iv) Echivalența modulo $\mathbb Q$ pe $\mathbb R$.}
Definim
\[
 x\sim y\overset{\mathrm{def}}\iff x-y\in\mathbb Q,
 \qquad \mathbb R/{\sim}\cong\mathbb R.
\]
Verificăm cele trei proprietăți:
\[
\begin{aligned}
\text{(R)}\quad &x\sim x\quad(0\in\mathbb R),\\
\text{(S)}\quad &x-y\in\mathbb Q\Longrightarrow y-x\in\mathbb Q,\\
\text{(T)}\quad &x-y\in\mathbb Q,\ y-z\in\mathbb Q
      \Longrightarrow x-z\in\mathbb Q.
\end{aligned}
\]
Proiecția este $p:\mathbb R\to\mathbb R/{\sim}$.
Notițele remarcă faptul că $\mathbb R$ nu este, ca atare,
un sistem de reprezentanți. Pentru un sistem de reprezentanți
$S=(x_i)_{i\in I}$,
\[
 p|_S:S\longrightarrow\mathbb R/{\sim}\quad\text{este bijecție}.
\]
Deoarece
\[
 \mathbb R=\bigcup_{i\in I}\widehat{x_i},\qquad
 \widehat{x_i}=\{y\in\mathbb R\mid y-x_i\in\mathbb Q\}
   =x_i+\mathbb Q\cong\mathbb Q,
\]
fiecare $\widehat{x_i}$ este o mulțime numărabilă. Dacă $I$ ar fi
numărabilă, atunci $\mathbb R$ ar fi numărabilă - contradicție.
Prin urmare, $I$ este nenumărabilă.

\endgroup
